\section{Anomalies présentes et lignes justificatives}\label{ann:qualite}
Fréquences issues de \chemin{results/quality/quality_report.json}, calculées
\emph{par fichier entier}, indépendamment des profils. Elles peuvent se recouper.
Nous avons comparé les huit exemples aux colonnes originales des Parquet et recalculé leurs durées avec New York puis UTC. Le relevé de cette vérification est
\chemin{report/evidence-verification.json}. Une transition DST est une particularité
à traiter, pas nécessairement une erreur de données. Les numéros de ligne sont en base zéro.

\begin{center}\small
\begin{tabular}{llrr}
\toprule
Mois & Indicateur & Fréquence & Part du fichier (\%)\\
\midrule
2024-01 & \code{PICKUP_OUTSIDE_FILE_MONTH} & 18 & 0,000607\\
2024-01 & \code{NON_POSITIVE_DURATION} & 870 & 0,029346\\
2024-01 & \code{ZERO_DISTANCE} & 60\,371 & 2,036380\\
2024-01 & \code{NEGATIVE_AMOUNT} & 35\,504 & 1,197589\\
2024-01 & \code{UNKNOWN_ZONE} & 31\,527 & 1,063440\\
2024-02 & \code{EXACT_DUPLICATE} & 1 & 0,000033\\
2024-03 & \code{DST_SPANNING} & 1\,153 & 0,032183\\
2024-03 & \code{DST_UNRESOLVED} & 3 & 0,000084\\
\bottomrule
\end{tabular}\end{center}

\paragraph{2024-01, ligne 256 --- \code{PICKUP_OUTSIDE_FILE_MONTH}}
Départ : 2023-12-31 23:56:46 ; arrivée : 2024-01-01 00:12:06. Zones 236 $\to$ 142, fournisseur 2, distance 2.38 miles, total 21.6.
Le départ est en décembre alors que la ligne est publiée en janvier ; la requête doit utiliser la date réelle.

\paragraph{2024-01, ligne 87\,994 --- \code{NON_POSITIVE_DURATION}}
Départ : 2024-01-02 10:00:00 ; arrivée : 2024-01-02 09:53:56. Zones 75 $\to$ 95, fournisseur 1, distance 7.2 miles, total 46.94.
L'arrivée précède le départ : durée réelle $-364$ s, exclue du profil journey.

\paragraph{2024-01, ligne 17 --- \code{ZERO_DISTANCE}}
Départ : 2024-01-01 00:52:09 ; arrivée : 2024-01-01 00:52:28. Zones 237 $\to$ 237, fournisseur 2, distance 0.0 miles, total 8.0.
Distance nulle malgré 19 s entre les deux heures : la ligne reste stockée mais n'alimente pas journey.

\paragraph{2024-01, ligne 99 --- \code{NEGATIVE_AMOUNT}}
Départ : 2024-01-01 00:18:24 ; arrivée : 2024-01-01 00:30:39. Zones 249 $\to$ 232, fournisseur 2, distance 2.16 miles, total -18.5.
Le montant total négatif est exclu des profils analytiques. Le rapport qualité repère aussi cette ligne comme écriture d'annulation.

\paragraph{2024-01, ligne 24 --- \code{UNKNOWN_ZONE}}
Départ : 2024-01-01 00:42:05 ; arrivée : 2024-01-01 01:16:49. Zones 263 $\to$ 265, fournisseur 1, distance 23.9 miles, total 127.94.
L'arrivée 265 ne désigne pas une zone précise de New York : exclusion du profil flows et donc journey.

\paragraph{2024-02, ligne 754\,252 --- \code{EXACT_DUPLICATE}}
Départ : 2024-02-08 18:23:50 ; arrivée : 2024-02-08 18:38:14. Zones 162 $\to$ 230, fournisseur 2, distance 0.81 miles, total 23.16.
Le rapport qualité signale cette ligne comme copie d'une ligne précédente du même fichier ; seule la copie est exclue. Les champs cités identifient l'exemple, le rapport JSON conserve toutes les colonnes.
La ligne identique antérieure se trouve à la position 754\,248.

\paragraph{2024-03, ligne 942\,762 --- \code{DST_SPANNING}}
Départ : 2024-03-09 23:38:43 ; arrivée : 2024-03-10 03:40:14. Zones 249 $\to$ 203, fournisseur 2, distance 31.35 miles, total 131.14.
Durée UTC de 10\,891 s contre 14\,491 s locales : une heure est retranchée. Le flag n'exclut pas automatiquement le trajet de journey.

\paragraph{2024-03, ligne 841\,170 --- \code{DST_UNRESOLVED}}
Départ : 2024-03-09 02:31:02 ; arrivée : 2024-03-10 02:09:54. Zones 114 $\to$ 48, fournisseur 2, distance 3.36 miles, total 27.72.
L'arrivée est à une heure inexistante le 10 mars. La durée exploitable est null et non les 85\,132 s de différence locale.
